\section{Relation to Prior Work}
\label{sec:related}

\paragraph{Datasets and representations.}
CircuitNet and OpenABC-D provide data for physical-design and synthesis
learning~\cite{chai2022circuitnet,chowdhury2021openabcd}. ForgeEDA supplies RTL, mapped netlists,
AIGs, and placed netlists~\cite{shi2025forgeeda}. AiEDA and CircuitOps provide flow-to-data and
graph-oriented infrastructure~\cite{qiu2025aieda,liang2023circuitops}; EDA-Schema-V2 organizes
stage-resolved physical data across technologies~\cite{shrestha2026edaschema}. The R2G benchmark
standardizes five graph views for representation comparisons~\cite{zhou2026r2g}. Our work focuses on
qualified sample production and evaluates a four-stage prediction product with explicit identity,
visibility, and label checks.

\paragraph{Flow agents and evolving benchmarks.}
ChatEDA and ORFS-agent study tool orchestration and parameter
optimization~\cite{he2024chateda,ghose2026orfsagent}. HighTide uses agent skills, incremental
RTL-to-GDS compilation, and persistent per-design rationale for evolving
curation~\cite{goldblatt2026hightide}. Tool-interactive RTL-to-GDS benchmarks further evaluate complete
EDA execution~\cite{deng2026rtlgds}. Our emphasis is the connection between validated repair adoption
and run-consistent graph admission, with separate measurements of acquisition yield, repair outcomes,
and semantic output quality.

\paragraph{Repair evaluation.}
PostEDA-Bench studies DRC fixing and PPA convergence with machine-checkable
outcomes~\cite{liu2026postedabench}. Trace2Skill evolves verifier-guided agent skills during task
execution~\cite{du2026trace2skill}, while EvoDRC evolves repair knowledge while addressing layout
violations~\cite{wu2026evodrc}. Our repair experiment examines a complementary setting: proposals are
learned and validated before deployment, then reused under a fixed action policy. Separating these
phases measures the transfer of executable repair knowledge and its deployment cost.
